\section{Contribution and Use to Other Researchers}
\label{sec:contribution}

The proposed design would contribute along three dimensions. First, geographic targeting: existing trade-flow evidence pools all EU partners, so corridor-specific leakage is diluted. By restricting the analysis to North Africa, the design isolates a corridor where regulatory asymmetry is large, transport costs are low, and production capacity in emission-intensive sectors already exists. A positive $\beta$ would indicate that the aggregate nulls documented by \citet{naegele2019} conceal heterogeneity, a finding directly relevant for calibrating CBAM border adjustments by partner country. A well-identified null would carry a different but equally informative implication: that free allocation has been sufficient to prevent trade diversion even where it is most likely.

Second, treatment measurement: by defining the treatment as $C_{st} = P_t \times (1 - a_{st})$, the design measures the carbon cost net of the allocation subsidy, rather than using a binary EU~ETS coverage indicator or the nominal EUA price. This distinction matters: during most of the sample, sectors on the carbon-leakage list received 100 per cent of the benchmark allocation for free, so their effective carbon cost was near zero despite a positive permit price. Studies that treat the nominal EUA price as the treatment overstate cost exposure for these sectors and understate the true treatment contrast across sectors. Researchers estimating competitiveness, innovation, or emissions effects in other contexts would benefit from the same adjustment.

Third, temporal coverage: the MSR-driven price acceleration from 2018 onward generated the strongest effective carbon cost in the EU~ETS's history, yet most existing studies end their samples before 2015. The proposed window through 2021 would test whether leakage effects emerge at higher price levels.

The triple-difference structure, with product-level exposure interacted with a corridor indicator and time-varying policy intensity, is adaptable to other bilateral trade questions where regulation differentially affects product categories and partners. The comparison-group choice, the free-allocation treatment, and the absorption logic are transferable to studies of CBAM effectiveness once the definitive regime generates post-treatment variation.
